\documentclass[11pt]{article}
\usepackage[margin=1in]{geometry}
\usepackage{amsmath,amssymb}
\usepackage[T1]{fontenc}
\usepackage{lmodern}
\usepackage{hyperref}

\title{Literature-implied partial answer to an LPR exponent question}
\author{Run fa\_banach\_001}
\date{2026-06-28}

\begin{document}
\maketitle
\sloppy

\paragraph{Status.}
This is a \textbf{literature-implied partial answer}, not a new proof packet
and not a full solution of exponent-independence for the vector-valued
Littlewood--Paley--Rubio de Francia property.

\paragraph{Source question.}
Hytönen--Torrea--Yakubovich, \emph{The Littlewood--Paley--Rubio de Francia
property of a Banach space for the case of equal intervals}
(arXiv:0807.2981), introduce the arbitrary-interval \(LPR_p\) property and
recall the known necessary conditions \(LPR_p \Rightarrow\) UMD and type \(2\).
In the introduction, lines 340--344 of the parsed source say that ``it is not
known whether there is an implication between \(LPR_p\) and \(LPR_q\)'' for
two different exponents \(p,q\in[2,\infty)\).
The paper then solves the equal-length interval variant, not the
arbitrary-interval exponent question.

\paragraph{Supporting theorem.}
Potapov--Sukochev--Xu, \emph{On the vector-valued Littlewood-Paley-Rubio de
Francia inequality} (arXiv:1104.2671), explicitly return to the same
Banach-space \(LPR_p\) property.  Their introduction cites the
Hytönen--Torrea--Yakubovich necessary condition and again records that it is
unknown whether \(LPR_p\) is independent of \(p\).  They then prove the upward
direction.  Theorem 3.1 states:
\[
  X\in LPR_q,\quad 2\le q<\infty
  \quad\Longrightarrow\quad
  X\in LPR_p\quad\text{for every }q\le p<\infty .
\]

\paragraph{Identification.}
This theorem gives a positive implication between two different \(LPR\)
exponents in the direction from a smaller exponent to every larger exponent.
It is therefore a literature answer to the upward half of the source's
two-exponent question.

\paragraph{Scope.}
The supporting theorem does not prove the reverse implication, so full
exponent-independence remains open here.  It also does not settle the broader
converse problem \(UMD+\)type \(2\Rightarrow LPR_p\) for arbitrary intervals.
The run contains a serious no-promotion attempt,
\path{attempts/1104.2671_lpr_umd_type2_bootstrap_barrier.md}, explaining
why the equal-length theorem does not formally bootstrap to arbitrary
intervals: the missing step is a dimension-free \(R\)-bounded summation over
dyadic scales.

\paragraph{Search evidence.}
The cheap run indexes were searched for arXiv:0807.2981, arXiv:1104.2671,
Littlewood--Paley--Rubio, Rubio de Francia, \(LPR_p\), \(LPR_q\), UMD, and
type \(2\).  The relevant existing run memory is the no-promotion attempt on
arXiv:1104.2671, not a durable packet for this partial literature answer.
External arXiv/web search located arXiv:1104.2671 as the decisive later
source and did not locate a full arbitrary-interval exponent-independence
solution.

\begin{thebibliography}{9}
\bibitem{HTY2008}
T. P. Hytönen, J. L. Torrea, and D. V. Yakubovich,
\emph{The Littlewood--Paley--Rubio de Francia property of a Banach space for
the case of equal intervals},
arXiv:0807.2981.

\bibitem{PSX2011}
D. Potapov, F. Sukochev, and Q. Xu,
\emph{On the vector-valued Littlewood-Paley-Rubio de Francia inequality},
arXiv:1104.2671.
\end{thebibliography}

\end{document}
